% =====================================================================
%  Bericht im Spezialisierungsmodul --- M.Sc. Data Analytics, JLU Gießen
%  Project JLU, Team Gamma (Projekt 29)
%
%  Bauen:   ./build.sh          (oder: latexmk -pdf main.tex)
%  Aufräumen: ./build.sh clean
%
%  Diese Datei ist die Wurzel. Inhalt steht in chapters/, Vorspann in
%  frontmatter/, Einstellungen in preamble/. Nichts hier hineinschreiben,
%  was Inhalt ist.
% =====================================================================
\documentclass[12pt,a4paper,ngerman]{article}

% Pakete -- bewusst schlank, damit der Bericht auch auf einer
% BasicTeX-Installation baut (siehe README.md).
%
% Zwei Engines, ein Quelltext: pdfTeX bekommt inputenc, XeTeX/LuaTeX nicht.
% fontenc[T1] gilt für beide, sonst fehlen den Schriften die Umlaute.

\usepackage{iftex}
\usepackage[T1]{fontenc}
\ifPDFTeX
  \usepackage[utf8]{inputenc}
\fi
\usepackage{lmodern}
\usepackage{microtype}

% XeTeX mit T1-Schriften setzt das direkte Zeichen ß als "SS".
% ��ber newunicodechar wird es auf \ss{} umgebogen -- wirkt auch unter pdfTeX.
\usepackage{newunicodechar}
\newunicodechar{ß}{\ss{}}

% Ohne diese Zeile bleibt die Sprache Englisch: die Klassenoption ngerman
% reicht bei article nicht, babel wird dadurch nicht geladen.
\usepackage[ngerman]{babel}

\usepackage[a4paper,left=3cm,right=2.5cm,top=2.5cm,bottom=2.5cm]{geometry}

\usepackage{graphicx}
\usepackage{booktabs}
\usepackage{longtable}
\usepackage{array}
\usepackage{caption}
\usepackage{enumitem}
\usepackage{amsmath}
\usepackage{xcolor}
\usepackage{csquotes}

\usepackage[
  backend=bibtex,
  style=authoryear,
  sorting=nyt,
  maxcitenames=2,
  maxbibnames=99,
  giveninits=true,
  uniquename=init,
  doi=true,
  url=true
]{biblatex}
\AtBeginBibliography{%
  \setlength{\bibitemsep}{2pt}%
  \setlength{\bibparsep}{0pt}%
}

\usepackage[hidelinks]{hyperref}

% Lange URLs und DOIs muessen umbrechen duerfen. Als unteilbare Box lief eine
% Adresse im Literaturverzeichnis 56 pt ueber den Rand (Bau vom 24.11.2026).
% biblatex erlaubt Umbrueche an Ziffern, Gross- und Kleinbuchstaben, ohne ein
% Trennzeichen zu setzen.
\DeclareFieldFormat{url}{URL\addcolon\space\url{#1}}
\setcounter{biburlnumpenalty}{100}
\setcounter{biburlucpenalty}{100}
\setcounter{biburllcpenalty}{100}

% Satz und Metadaten

\setlength{\parindent}{0pt}
\setlength{\parskip}{0.55em}
\setlength{\emergencystretch}{2em}

\captionsetup{font=small,labelfont=bf}
\setlist{itemsep=0.15em,topsep=0.25em}
\renewcommand*{\bibfont}{\small}

% Zeilenumbruch in Tabellen nicht erzwingen
\renewcommand{\arraystretch}{1.15}

\hypersetup{
  pdftitle={Konsistenz über Quellen},
  pdfauthor={The Trung Do},
  pdfsubject={Spezialisierungsmodul M.Sc. Data Analytics, Project JLU, Team Gamma},
  pdfkeywords={Informationsqualität, Konsistenz, Abweichungserkennung, Studieninformationen}
}

% Helfer

% Arbeitsmarkierung. Vor der Abgabe muss `grep -r "\\todo" .` leer sein.
\newcommand{\todo}[1]{\textcolor{red}{\textbf{[TODO:} #1\textbf{]}}}

% Arbeitstitel, einmal zentral --- Titelblatt und Metadaten greifen darauf zu.
\newcommand{\arbeitstitel}{Konsistenz über Quellen}

% Sachverhalt im Fließtext (Attribute, Werte)
\newcommand{\fakt}[1]{\texttt{#1}}

% Abweichungstyp kurz referenzieren
\newcommand{\abwtyp}[1]{\textsc{#1}}


\addbibresource{literatur.bib}

\begin{document}

% ---------- Vorspann -------------------------------------------------
% Die Titelseite bekommt keinen Seitenanker: sie trägt die Nummer 1, und
% Kapitel 1 trägt sie später noch einmal -- hyperref würde "page.1" doppelt
% definieren. Kanonischer Weg laut hyperref-Handbuch.
\hypersetup{pageanchor=false}
% Titelblatt
% Der Titel ist ein Arbeitstitel. Nach der Abstimmung mit der Betreuung
% hier und in preamble/commands.tex (\arbeitstitel) anpassen.

\begin{titlepage}
  \centering
  \vspace*{1.2cm}

  {\scshape\large Justus-Liebig-Universität Gießen\par}
  \vspace{0.3cm}
  {\scshape\small Fachbereich 07 --- Mathematik und Informatik, Physik und Geographie\par}

  \vspace{2.6cm}

  {\huge\bfseries \arbeitstitel\par}
  \vspace{0.7cm}
  {\large Abweichungen zwischen Studieninformationen automatisch finden und bewerten\par}
  \vspace{0.5cm}
  {\small\textit{Arbeitstitel --- endgültige Formulierung nach Abstimmung mit der Betreuung}\par}

  \vspace{2.4cm}

  {\large Bericht im Spezialisierungsmodul\par}
  \vspace{0.2cm}
  {\normalsize M.Sc. Data Analytics, 12 Credit Points\par}

  \vspace{1.8cm}

  \begin{tabular}{@{}rl@{}}
    Vorgelegt von:   & The Trung Do \\
    Matrikelnummer:  & \matrikelnummer \\
    Betreuung:       & Dr. Sebastian Busse \\
    Projekt:         & Project JLU, Team Gamma (Projekt 29) \\
    Abgabedatum:     & 1. April 2027 \\
  \end{tabular}

  \vfill
  {\small Gießen, \today\par}
\end{titlepage}

\hypersetup{pageanchor=true}

% Römische Zählung für den Vorspann, arabische für den Text.
\pagenumbering{roman}
% Zusammenfassung --- 120 bis 180 Wörter.
% Erst schreiben, wenn die Ergebnisse stehen (März), nicht vorher.

\section*{Zusammenfassung}
\addcontentsline{toc}{section}{Zusammenfassung}

\todo{Zusammenfassung: (1) Fragestellung in einem Satz, (2) Datengrundlage und Stichtag,
(3) Verfahren in zwei Sätzen, (4) die wichtigste Zahl mit Unsicherheit, (5) Einordnung in
einem Satz. Keine Wertungen, keine Ausblicke ohne Befund.}

\clearpage

\tableofcontents
\clearpage

% ---------- Hauptteil ------------------------------------------------
\pagenumbering{arabic}
% Kapitel 1 --- Einleitung
% Budget: 1,5 Seiten (von 13 ohne Literatur und Anhang)
% Quelle im Vault: 10 Thema/Thema 29 -- Überblick, 10 Thema/3.1 Konsistenz

\section{Einleitung}
\label{sec:einleitung}

\subsection{Problem}

Studieninformationen stehen nicht an einer Stelle. Dieselbe Angabe --- eine Bewerbungsfrist,
eine Zulassungsvoraussetzung, die Zahl der Leistungspunkte eines Moduls --- erscheint in der
Prüfungsordnung, auf der Studiengangsseite und in den Informationsschriften der zentralen
Studienberatung. Die Dokumente entstehen zu verschiedenen Zeiten, werden von verschiedenen
Stellen gepflegt und nur selten gemeinsam fortgeschrieben. Für die Leserin ist das
unsichtbar: Jede Seite wirkt für sich vollständig und aktuell, und es gibt keine Stelle, die
die Angaben gegeneinander hält.

Ein Beispiel aus dem Korpus dieser Arbeit zeigt, wie unauffällig solche Abweichungen
entstehen.\footnote{Übersichtsseite der Mitteilungen der JLU 7.36.07 Nr.~11 und PDF der
1.~Änderungsfassung; die Adressen sind in Abschnitt~\ref{sec:daten} verzeichnet.} Die
Spezielle Ordnung des Masterstudiengangs Data Analytics liegt in vier Fassungen vor. Für die
Fassung des 1.~Änderungsbeschlusses nennt die Übersichtsseite der Mitteilungen den
Geltungsbeginn \enquote{gültig ab dem Wintersemester 2023/24}; die Titelseite derselben
Fassung formuliert \enquote{gilt ab dem Wintersemester 2024/25}. Beide Angaben stehen auf
Servern derselben Hochschule, beide sind zitierfähig, und beide können nicht zutreffen. Wer
die Übersichtsseite liest, hält die Änderung ein Jahr früher für verbindlich als jemand, der
die Ordnung selbst aufschlägt.

Die Ursache ist strukturell und nicht persönlich: Die Übersichtsseite fasst zusammen, die
Ordnung ist der verbindliche Text, und beide werden von Hand gepflegt. Solche Stellen sind
der Gegenstand dieser Arbeit. Sie fragt nicht, ob die Universität ihre Seiten gut pflegt,
sondern ob sich Abweichungen dieses Typs systematisch finden, zuverlässig bewerten und nach
ihrer Entstehung unterscheiden lassen.

\subsection{Relevanz}

Informationsqualität ist keine Geschmacksfrage. ISO/IEC 25012 führt \emph{Consistency} als
eigenes Qualitätsmerkmal: Daten sind konsistent, wenn sie frei von Widersprüchen und mit
anderen Daten kohärent sind \parencite{iso25012}. Ein Datenprodukt, das sich selbst
widerspricht, ist damit unabhängig davon mangelhaft, ob jemand den Widerspruch bemerkt.

Wie schnell solche Widersprüche entstehen, ist für wissenschaftliche Verweise gut
dokumentiert. \textcite{klein2014} finden Reference Rot in jedem fünften der von ihnen
untersuchten Artikel; \textcite{jones2016} berichten, dass drei von vier URI-Verweisen auf
Inhalte zeigen, die sich gegenüber dem Zitat verändert haben. Die Studien betreffen
Literaturverweise und nicht Studieninformationen, und ihre Zahlen lassen sich nicht
übertragen. Sie belegen aber den Mechanismus: Dokumente werden überarbeitet, verschoben und
ersetzt, während die Verweise auf sie stehen bleiben.

Für eine Hochschule kommt ein zweiter Punkt hinzu. Eine Studienentscheidung ist folgenreich
und wird in der Regel einmal getroffen: Wer sich auf eine überholte Fassung stützt, merkt es
nicht an der Seite, sondern frühestens im Verfahren. Der Aufwand, solche Stellen zu finden,
ist heute Handarbeit --- und er wächst mit der Zahl der Dokumente. Genau dort setzt die
Automatisierung an, die diese Arbeit untersucht.

\subsection{Beitrag dieser Arbeit}

Die Arbeit leistet drei Beiträge. Erstens ein Verfahren, das Angaben aus verschiedenen
Dokumenten extrahiert, normalisiert und nach Sachverhalten gruppiert, sodass Abweichungen
systematisch sichtbar werden statt zufällig aufzufallen. Zweitens eine Messung der
Zuverlässigkeit dieses Verfahrens: Gegen einen von Hand erstellten Goldstandard werden
Precision und Recall je Sachverhaltstyp sowie die Übereinstimmung zwischen Mensch und Modell
berichtet. Drittens eine Typologie der Abweichung, die echte Widersprüche von
Versionsartefakten, Scope-Unterschieden und Extraktionsfehlern trennt --- erst sie macht aus
einer Fundliste einen Befund.

\subsection{Aufbau}

Kapitel~\ref{sec:hintergrund} ordnet die Arbeit in Informationsqualität,
Widerspruchserkennung und Content Drift ein und grenzt sie von den Nachbardimensionen ab.
Kapitel~\ref{sec:frage} formuliert Forschungsfrage und Hypothesen samt
Widerlegungskriterien. Kapitel~\ref{sec:daten} beschreibt Korpus, Faktkatalog und
Erhebungsgrenzen, Kapitel~\ref{sec:methode} die Analyseeinheit, die Extraktion mit
Zitatpflicht, die Normalisierung, die Abweichungstypologie und das Evaluationsdesign.
Kapitel~\ref{sec:ergebnisse} berichtet die Messung, Kapitel~\ref{sec:diskussion} ordnet sie
ein und benennt die Grenzen, Kapitel~\ref{sec:fazit} fasst das Ergebnis zusammen.

% Kapitel 2 --- Hintergrund und verwandte Arbeiten
% Budget: 2,5 Seiten

\section{Hintergrund und verwandte Arbeiten}
\label{sec:hintergrund}

\subsection{Informationsqualität als Rahmen}

Informationsqualität ist kein einzelnes Merkmal, sondern ein Bündel von Dimensionen.
\textcite{wang1996beyond} leiten aus Befragungen fünfzehn Dimensionen ab und ordnen sie in
intrinsische, kontextuelle, repräsentationale und zugangsbezogene Kategorien. Der
Standard ISO/IEC 25012 \parencite{iso25012} fasst dieselbe Idee normativ und definiert
unter anderem \emph{Consistency} als Freiheit von Widerspruch und Kohärenz mit anderen
Daten sowie \emph{Currentness} als richtiges Alter im Nutzungskontext. ISO/IEC 25024
\parencite{iso25024} ergänzt Messgrößen zu diesen Merkmalen.

Das Projekt untersucht fünf dieser Dimensionen getrennt: Konsistenz, Aktualität,
Auffindbarkeit, Verständlichkeit und Mehrsprachigkeit. Die vorliegende Arbeit bearbeitet die
Konsistenz, aus drei Gründen. Sie ist normativ verankert und braucht keine eigene Definition
des Qualitätsbegriffs; sie hat eine natürliche Analyseeinheit, den Sachverhalt, der sich
zählen lässt; und sie ist quellenübergreifend, lässt sich also nicht durch die Prüfung eines
einzelnen Dokuments erledigen. Wichtig ist dabei die Arbeitsteilung: Der Standard sagt,
\emph{was} gemessen wird, nicht \emph{wie}. Die Operationalisierung --- Extraktion,
Normalisierung, Typologie, Goldstandard --- ist der Beitrag dieser Arbeit und nicht aus der
Norm ableitbar.

\subsection{Widersprüche automatisch erkennen}

Der Referenzrahmen für automatische Widerspruchsprüfung ist die Aufgabenstellung von FEVER
\parencite{thorne2018fever}: Zu einer Behauptung wird eine Belegstelle gesucht und daraus ein
Urteil abgeleitet --- gestützt, widerlegt oder nicht entscheidbar. Der Datensatz umfasst
185.445 Behauptungen, die gegen Absätze aus Wikipedia geprüft werden. Für diese Arbeit ist
nicht der Umfang entscheidend, sondern die Struktur: Das Urteil hängt an einer zitierbaren
Stelle und nicht an der Überzeugung des Prüfers. Genau diese Struktur braucht die
Konsistenzmessung, denn sie vergleicht nicht eine Aussage mit Weltwissen, sondern zwei
Dokumente miteinander.

Für Fließtext, in dem Werte nicht exakt übereinstimmen, ist Natural Language Inference ein
Baustein. SNLI stellt rund 570.000 Satzpaare mit den Urteilen \emph{Implikation}, \emph{neutrale
Aussage} und \emph{Widerspruch} bereit \parencite{bowman2015snli}; XNLI erweitert das
Verfahren auf fünfzehn Sprachen \parencite{conneau2018xnli}. Der Nutzen für diese Arbeit ist
begrenzt und klar umrissen: NLI entscheidet, ob zwei Sätze einander widersprechen, wenn sie
verschieden formuliert sind. Die Prüfung, ob sie denselben Sachverhalt betreffen, ersetzt es
nicht.

Die Grenzen sind empirisch belegt. \textcite{pielka2022} zeigen, dass Widerspruchsmodelle an
Präpositionen und Verben, an Antonymen und Homonymen, an unvollständigen Sätzen, an langen
Absätzen und an seltenen Wörtern scheitern, und dass sich crowdsourced-maschinell übersetzte
Daten sprachlich von Textpaaren aus Internetquellen unterscheiden. Daraus folgt für diese
Arbeit zweierlei: Der Goldstandard muss aus dem eigenen Korpus stammen, und ein
SNLI-trainiertes Urteil ist kein Beleg für einen Widerspruch zwischen zwei JLU-Seiten. Es ist
ein Hinweis, der an der Fundstelle geprüft werden muss.

\subsection{Veraltung und Content Drift}

Zwei Phänomene werden häufig vermischt. \emph{Link Rot} bezeichnet Verweise, die ins Leere
führen; \emph{Content Drift} bezeichnet Verweise, die weiter funktionieren, aber auf
veränderten Inhalt zeigen \parencite{klein2014,jones2016}. Für diese Arbeit ist der zweite
Fall der relevante: Eine Seite bleibt erreichbar und sieht gepflegt aus, während die Angabe
darauf nicht mehr stimmt.

Wie ungleich solche Änderungen verteilt sind, zeigt \textcite{fetterly2003} an 150.836.209
HTML-Seiten über elf Wochen: Die Änderungsrate variiert stark zwischen Top-Level-Domains, und
größere Seiten ändern sich häufiger als kleinere. \textcite{cho2000freshness} zeigen
anschließend, wie sich Abrufpläne an die gemessene Änderungsrate anpassen lassen. Übertragen
auf Studieninformationen heißt das: Nicht jede Angabe veraltet gleich schnell, und ein
Vergleich über alle Angaben hinweg mischt schnelle und langsame Typen. Ob Fristen tatsächlich
häufiger abweichen als Modulnummern, ist eine der Hypothesen dieser Arbeit
(Abschnitt~\ref{sec:frage}); die Literatur zeigt nur, dass solche Unterschiede zu erwarten
sind.

Für die Fragestellung folgt daraus eine Unterscheidung, die die ganze Auswertung trägt: Ein
Teil der gefundenen Abweichungen ist kein Widerspruch, sondern eine nicht mitgezogene
Aktualisierung. Beide Fälle sehen im Datensatz gleich aus --- zwei verschiedene Werte für
denselben Sachverhalt --- und müssen dennoch verschieden behandelt werden. Die Typologie in
Abschnitt~\ref{sec:methode} trennt sie.

\subsection{Abgrenzung}

Drei Abgrenzungen sind wichtig. Verweise und Fristen \emph{innerhalb} der Ordnungen --- ob
eine Ordnung auf eine andere verweist und dabei die richtige Fassung nennt --- prüft die
Arbeit 2.2 desselben Projekts; hier geht es ausschließlich um Abweichungen \emph{zwischen}
Dokumenten. Ob zwei Erwähnungen dieselbe Entität meinen, ist Gegenstand von 2.4 (Entity
Linking); diese Arbeit setzt die Entitätsauflösung über eine gepflegte Aliasliste voraus,
statt sie zu bewerten. Und die Veraltung als eigene Dimension mit Zeitreihe und Vorhersage
ist 3.2 --- hier ist sie nur ein Erklärungsmuster für einen Teil der Abweichungen.

% Kapitel 3 --- Forschungsfrage und Hypothesen
% Budget: 1 Seite
% Quelle im Vault: 20 Recherche/Hypothesen 3.1

\section{Forschungsfrage und Hypothesen}
\label{sec:frage}

\subsection{Fragestellung}

\todo{Die Fragestellung aus dem Steckbrief aufnehmen und auf die eigene Operationalisierung
zuspitzen: Welche Sachverhalte werden in mehreren Quellen unterschiedlich dargestellt, und
wie zuverlässig lässt sich das automatisch finden?}

\subsection{Hypothesen}

\begin{description}[style=nextline]
  \item[H1 --- Abweichungen konzentrieren sich auf instabile Sachverhaltstypen]
  \todo{Formulierung, Begründung aus der Literatur, Operationalisierung, und was als
  Widerlegung gilt.}
  \item[H2 --- Sprachmodellgestützte Extraktion ist präzise, aber nicht vollständig]
  \todo{Formulierung, Bezug zur Validierungsliteratur \parencite{bavaresco2025llms},
  Operationalisierung über Precision und Recall je Verfahrensvariante.}
  \item[H3 --- Die meisten Abweichungen sind Zeitartefakte]
  \todo{Formulierung, Begründung über Content Drift und die Änderungsfassungen der
  Ordnung, Operationalisierung über die Abweichungstypologie.}
\end{description}

\subsection{Was als Widerlegung gilt}

\todo{Für jede Hypothese einen Satz: Welches Ergebnis würde sie kippen? Ohne diesen Absatz
ist die Arbeit nicht falsifizierbar, und das fällt in der Bewertung auf.}

% Kapitel 4 --- Daten und Korpus
% Budget: 1,5 Seiten
% Quelle im Vault: 20 Recherche/Datenquellen, 10 Thema/Team Gamma -- gemeinsame Grundlage

\section{Daten und Korpus}
\label{sec:daten}

\subsection{Quellen und Stichtag}

\todo{Welche Dokumente, welche Fassungen, ein Abrufdatum. Tabelle mit Quelle, Typ, Anzahl,
Abrufdatum. Die Fassungen der Speziellen Ordnung einzeln nennen --- sie sind der Zeitstrahl
der Arbeit.}

\subsection{Faktkatalog}

\todo{Die untersuchten Sachverhaltstypen mit Begründung der Auswahl. Vollständige Tabelle
in den Anhang, hier nur die Logik: warum diese zwölf und nicht vierzig.}

\subsection{Ausschlüsse und Grenzen der Erhebung}

\todo{Was nicht erfasst wurde: Seiten hinter Login, PDFs ohne Textebene, dynamische
Inhalte. Ehrlich und vollständig --- dieser Absatz entscheidet mit über die Bewertung
der Belastbarkeit.}

\subsection{Erhebung und Höflichkeitsregeln}

\todo{robots.txt, Rate, User-Agent, Ankündigung bei der Webredaktion. Kurz, aber
nachweisbar.}

% Kapitel 5 --- Methode
% Budget: 3 Seiten --- das Herzstück, hier entscheidet sich die Note
% Quelle im Vault: 20 Recherche/Methodik

\section{Methode}
\label{sec:methode}

\subsection{Analyseeinheit: der Sachverhalt}

\todo{Definition als Tupel aus Entität, Attribut, Wert, Gültigkeitsintervall, Quelle und
Fundstelle. Ein Beispiel ausgeschrieben. Begründung: Ohne diese Definition bleibt
\enquote{derselbe Sachverhalt} Alltagssprache und die Messung wird nicht prüfbar.}

\subsection{Extraktion mit Zitatpflicht}

\todo{Sprachmodell mit festem Schema, jeder Wert mit Belegstelle; Werte ohne Fundstelle
gelten als nicht extrahiert. Begründung: Ein Modell liefert plausible Werte auch dort,
wo keine stehen. Temperatur, Prompt-Version, Modellversion --- gehört in den Anhang.}

\subsection{Normalisierung}

\todo{Datum, Zahl, Einheit, Entitätsalias, Scope. Tabelle der Regeln mit je einem
Beispiel. Hinweis: Die Aliasliste ist selbst ein Ergebnis, weil sie zeigt, wie viele
Schreibweisen derselbe Studiengang hat.}

\subsection{Abweichungstypologie}

\todo{Die fünf Typen (echter Widerspruch, Versionsunterschied, Scope-Unterschied,
Formulierung, Extraktionsfehler) mit Entscheidungsregel je Typ. Das ist der eigene
Beitrag --- hier ausführlich werden.}

\subsection{Goldstandard und Reliabilität}

\todo{Stichprobe, zwei Durchgänge, Konfliktlösung, Krippendorffs Alpha mit Konvention
\parencite{krippendorff2004}. Und die Validierung des Modells gegen menschliche Labels
\parencite{bavaresco2025llms}.}

\subsection{Metriken}

\todo{Precision, Recall, F1 je Sachverhaltstyp; Wilson-Konfidenzintervall; Formeln
hinschreiben, nicht nur nennen. Warum je Typ und nicht global.}

% Kapitel 6 --- Ergebnisse
% Budget: 3 Seiten
% Regel: Jede Zahl kommt aus einem Skriptlauf, nicht aus dem Kopf.

\section{Ergebnisse}
\label{sec:ergebnisse}

\subsection{Abweichungsraten je Sachverhaltstyp}

\todo{Tabelle mit Typ, Anzahl Sachverhalte, Abweichungen, Anteil, 95-Prozent-Intervall.
Ein Satz je Zeile, der die Zahl einordnet. Die Tabelle ist das Ergebnis, nicht die
Vorbereitung darauf.}

\subsection{Verfahrensvergleich}

\todo{Nur Sprachmodell, nur Regeln, Hybrid: Precision, Recall, F1 gegen denselben
Goldstandard. Aussage: Was bringt die Normalisierung wirklich?}

\subsection{Fehleranalyse}

\todo{20 bis 30 Fehlfälle, kategorisiert und gezählt. Diese Tabelle zeigt, dass das
eigene Verfahren verstanden wird --- sie ist oft der stärkste Abschnitt einer Arbeit.}

\subsection{Beispiele}

\todo{Zwei Fälle ausgeschrieben: ein echter Widerspruch und ein Versionsartefakt, jeweils
mit Fundstellen und Typzuordnung. Nicht mehr als zwei.}

% Kapitel 7 --- Diskussion
% Budget: 1,5 Seiten

\section{Diskussion}
\label{sec:diskussion}

\subsection{Interpretation}

\todo{Was bedeuten die Zahlen für die JLU? Ein Absatz, der ohne die Zahlen davor nicht
verständlich wäre --- also nicht wiederholen, sondern einordnen.}

\subsection{Prüfung der Hypothesen}

\todo{Je Hypothese: bestätigt, teilweise bestätigt, widerlegt --- mit der Zahl, die
entscheidet, und ohne Schönfärberei.}

\subsection{Grenzen}

\todo{Fünf ehrliche Punkte: Stichprobengröße, Auswahl der Sachverhalte, Modell als
Messinstrument, ein Stichtag statt Zeitreihe, keine Nutzerperspektive. Wer die Grenzen
selbst benennt, wird nicht darauf festgenagelt.}

\subsection{Anschlussfähigkeit}

\todo{Was Team Beta damit tun könnte (Entity Linking, Graphqualität), und was eine
Folgearbeit bräuchte. Zwei Sätze, kein Ausblick-Kapitel.}

% Kapitel 8 --- Fazit
% Budget: 0,5 Seiten

\section{Fazit}
\label{sec:fazit}

\todo{Drei Absätze, je zwei bis drei Sätze: (1) Was wurde untersucht und wie. (2) Was kam
heraus --- die eine Zahl, die man sich merken soll. (3) Was folgt daraus für die JLU.
Keine Wiederholung der Einleitung, keine neuen Behauptungen.}


% ---------- Nachspann ------------------------------------------------
\clearpage
% Anhang --- zählt nicht zu den 10 bis 15 Seiten.
% Alles, was Beleg ist und nicht Argument, gehört hierher.

\appendix
\section{Anhang}

\subsection{Faktkatalog}
\label{app:faktkatalog}
\todo{Tabelle der untersuchten Sachverhaltstypen: Typ, Beispiel, erwartete
Änderungsrate, Quellen. Aus dem Team-Repo übernehmen, sobald es steht.}

\subsection{Extraktionsschema}
\label{app:schema}

Das Schema ist die technische Form des Sachverhalts-Tupels aus
Abschnitt~\ref{sec:methode}. Es hat drei Ebenen: das Dokument, den Fakt und die Fundstelle.
Tabelle~\ref{tab:schema} nennt die Felder.

\begin{table}[ht]
  \centering
  \caption{Felder des Extraktionsschemas}
  \label{tab:schema}
  \small
  \begin{tabular}{@{}>{\raggedright\arraybackslash}p{0.24\textwidth}>{\raggedright\arraybackslash}p{0.22\textwidth}c>{\raggedright\arraybackslash}p{0.34\textwidth}@{}}
    \toprule
    Feld & Typ & Pflicht & Regel \\
    \midrule
    doc\_id & Zeichenkette & ja & Kennung aus dem Datenblatt, z.\,B.\ \texttt{ord\_3ae} \\
    doc\_hash & SHA-256 & ja & bindet den Fakt an die eingefrorene Fassung \\
    entitaet & Zeichenkette & ja & kanonisch nach Aliasliste \\
    attribut & kontrolliertes Vokabular & ja & z.\,B.\ \texttt{bewerbungsfrist\_wise} \\
    wert\_raw & Zeichenkette & ja & wörtlich, wie im Dokument \\
    wert\_norm & Datum, Zahl, Text & nein & Ergebnis der Normalisierung \\
    einheit & Zeichenkette & nein & LP, Minuten, EUR \\
    scope & Zeichenkette & ja & Zielgruppe oder Programm; Vorgabe \texttt{alle} \\
    gueltig\_von & Datum & nein & Beginn des Gültigkeitsintervalls \\
    gueltig\_bis & Datum & nein & Ende des Gültigkeitsintervalls \\
    fundstelle\_text & Zeichenkette & ja & wörtliches Zitat \\
    fundstelle\_start & ganze Zahl & ja & Zeichenoffset im segmentierten Text \\
    fundstelle\_ende & ganze Zahl & ja & Zeichenoffset im segmentierten Text \\
    status & Aufzählung & ja & \texttt{ok} oder \texttt{verworfen} \\
    prompt\_version & Zeichenkette & ja & Reproduzierbarkeit des Laufs \\
    modell & Zeichenkette & ja & Name und Version des Modells \\
    \bottomrule
  \end{tabular}
\end{table}

Die Pflichtfelder sind so geschnitten, dass jeder Fakt ohne Rückfrage prüfbar ist:
\texttt{doc\_hash} und die beiden Offsets verweisen auf eine Stelle in einer eingefrorenen
Datei, nicht auf ein Dokument, das sich inzwischen geändert haben kann. \texttt{wert\_norm}
ist bewusst nicht Pflicht. Die Normalisierung darf scheitern, ohne den Fakt zu verlieren:
Der Rohwert bleibt erhalten, der Fall wird als nicht normalisierbar markiert und in der
Stichprobe manuell geprüft. Ein Fakt ohne Fundstelle wird dagegen nicht geschrieben --- er
gilt als nicht extrahiert.

Die Prüfung erfolgt beim Schreiben. Ein Datensatz, der ein Pflichtfeld nicht erfüllt, wird
verworfen und protokolliert, nicht korrigiert; damit ist die Datenbank eine Teilmenge des
tatsächlich Gelesenen und keine Interpretation. \texttt{wert\_raw} und \texttt{wert\_norm}
stehen nebeneinander, sodass sich jede Regel aus Abschnitt~\ref{sec:methode} am Einzelfall
nachvollziehen lässt; wird eine Regel geändert, entstehen neue normalisierte Werte, ohne
dass die Extraktion wiederholt werden muss. Der Hash schließlich trennt Dokumentversionen:
Ein zweiter Abruf derselben Adresse ist eine neue Fassung und keine Korrektur der alten ---
genau die Unterscheidung, die für Versionsartefakte gebraucht wird.

\subsection{Verwendete Sprachmodelle und Prompts}
\label{app:modelle}
\todo{Modellname, Version, Temperatur, Prompt-Version, Datum. Pflicht, wenn ein
Sprachmodell im Spiel ist --- sonst ist die Arbeit nicht reproduzierbar.}

\subsection{Fehlertaxonomie}
\label{app:fehlertaxonomie}

Die Taxonomie beschreibt Fehler des Verfahrens, nicht Abweichungen der Quellen. Sie ist von
der Typologie in Abschnitt~\ref{sec:typologie} zu unterscheiden: Dort geht es darum, was ein
Fund bedeutet; hier darum, wodurch das Verfahren danebenliegt. Tabelle~\ref{tab:fehler}
nennt die fünf Klassen.

\begin{table}[ht]
  \centering
  \caption{Fehlerklassen des Verfahrens}
  \label{tab:fehler}
  \small
  \begin{tabular}{@{}>{\raggedright\arraybackslash}p{0.20\textwidth}>{\raggedright\arraybackslash}p{0.44\textwidth}>{\raggedright\arraybackslash}p{0.26\textwidth}@{}}
    \toprule
    Klasse & Definition & Folge für die Kennzahl \\
    \midrule
    Extraktion & Der Wert steht so nicht im Dokument; das Zitat trägt ihn nicht & falscher Widerspruch (Precision) oder fehlender Fakt (Recall) \\
    Erfundene Angabe & Wert ohne Deckung im Text; trotz Zitatpflicht möglich, wenn das Zitat nicht zum Wert passt & Precision \\
    Normalisierung & richtig gelesen, falsch vereinheitlicht, etwa 1,5 Stunden als 90 statt 150 Minuten & Precision (falscher Widerspruch) \\
    Gruppierung & verschiedene Sachverhalte zusammengeworfen oder derselbe getrennt & Precision und Recall \\
    Scope & Zielgruppe übersehen, etwa EU-Bewerbende gegen alle & Precision \\
    Urteil & Widerspruch gesehen, wo keiner ist, oder einen übersehen & Precision bzw.\ Recall \\
    \bottomrule
  \end{tabular}
\end{table}

Die Klassen sind vor der Auswertung festgelegt, nicht danach. Der Grund ist methodisch: Wer
die Kategorien erst nach dem Blick auf die Fehlfälle bildet, passt sie unmerklich dem
Ergebnis an, und die Fehleranalyse verliert ihre Funktion als Kontrolle. Jeder Fall bekommt
eine Hauptursache, weil sonst ein einziger Fehler mehrfach gezählt würde; Nebenursachen
werden im Protokoll vermerkt.

Die Abgrenzung zur Abweichungstypologie ist wichtig. Ein Extraktionsfehler ist dort der
erste Typ und führt zum Ausschluss des Falls; in der Fehlertaxonomie ist er eine Messung der
Verfahrensqualität. Dieselbe Beobachtung kann also an zwei Stellen auftauchen, aber mit
verschiedener Funktion: einmal als Grund, einen Fund zu verwerfen, einmal als Kennzahl
dafür, wie oft das nötig war. Die Häufigkeit je Klasse wird erst nach dem Lauf berichtet;
sie ist eine Messzahl und steht deshalb in Kapitel~\ref{sec:ergebnisse} und nicht hier.

\subsection{Code-Verzeichnis}
\label{app:code}

Der Bericht wird aus dem Verzeichnis \texttt{report/} gesetzt. Die Auswertungswerkzeuge
entstehen in den Arbeitspaketen P2 und P3; die Verzeichnisse für Korpus, Code und Tests sind
noch nicht angelegt. Tabelle~\ref{tab:dateien} verzeichnet deshalb den heutigen Bestand und
keinen geplanten.

\begin{table}[ht]
  \centering
  \caption{Dateien und Verzeichnisse des Berichts}
  \label{tab:dateien}
  \small
  \begin{tabular}{@{}>{\raggedright\arraybackslash}p{0.30\textwidth}>{\raggedright\arraybackslash}p{0.62\textwidth}@{}}
    \toprule
    Datei oder Verzeichnis & Rolle \\
    \midrule
    \texttt{build.sh} & baut \texttt{main.pdf}; Ziele \texttt{expose} und \texttt{clean} \\
    \texttt{main.tex} & Wurzel; bindet Vorspann, Kapitel, Anhang und Literatur ein \\
    \texttt{literatur.bib} & Literaturangaben; 32 Einträge, davon 13 im Bericht zitiert \\
    \texttt{preamble/} & Pakete, Satzspiegel und Kommandos (drei Dateien) \\
    \texttt{frontmatter/} & Titelseite, Zusammenfassung, Selbstständigkeitserklärung \\
    \texttt{chapters/} & acht Kapiteldateien, \texttt{01\_einleitung} bis \texttt{08\_fazit} \\
    \texttt{backmatter/} & Anhang A.1 bis A.5 \\
    \texttt{figures/}, \texttt{tables/} & Ablage für Abbildungen und Tabellen; derzeit leer \\
    \texttt{expose.tex} & Ein-Seiten-Exposé als eigenes Dokument mit derselben Präambel \\
    \texttt{README.md} & Bauanleitung und Hinweise zur Engine \\
    \texttt{.gitignore} & Build-Artefakte; \texttt{main.pdf} wird mitgeführt \\
    \bottomrule
  \end{tabular}
\end{table}

Ein Code-Verzeichnis im engeren Sinn ist dieser Anhang noch nicht. Solange die Pipeline nicht
läuft, wäre eine Liste von Modulen eine Behauptung über Dateien, die es nicht gibt. Sobald
die Werkzeuge stehen, wird die Tabelle um ihre Einträge ergänzt; die Rollen der vorhandenen
Dateien bleiben dabei unverändert.


\clearpage
\printbibliography[heading=bibintoc,title={Literaturverzeichnis}]

\clearpage
% Selbstständigkeitserklärung
%
% ACHTUNG: Der Wortlaut muss aus der Vorlage des Moduls übernommen werden.
% Der folgende Text ist ein Platzhalter und **nicht** die geprüfte Fassung.
% Vor der Abgabe: Vorlage der Betreuung anfordern, Text ersetzen, nichts ergänzen,
% was dort nicht steht.

\section*{Selbstständigkeitserklärung}
\addcontentsline{toc}{section}{Selbstständigkeitserklärung}

Ich versichere, dass ich die vorliegende Arbeit selbstständig verfasst und keine anderen
als die angegebenen Quellen und Hilfsmittel verwendet habe. Alle Stellen, die wörtlich
oder sinngemäß aus Veröffentlichungen entnommen sind, sind als solche kenntlich gemacht.

\vspace{0.6cm}
\noindent
Die verwendeten Sprachmodelle und Werkzeuge sind im Anhang aufgeführt; die Konzeption der
Arbeit, die Auswahl und Prüfung der Daten sowie die Interpretation der Ergebnisse stammen
von mir.

\vspace{2.2cm}
\noindent
\begin{minipage}{0.45\textwidth}
  \rule{\linewidth}{0.4pt}\\
  \small Ort, Datum
\end{minipage}
\hfill
\begin{minipage}{0.45\textwidth}
  \rule{\linewidth}{0.4pt}\\
  \small The Trung Do
\end{minipage}


\end{document}
